\subsection{Reproducibility audit of the internal benchmark}
The stratified 25\% holdout contains 494 patients and 200 events in each of the five seeds. Table~\ref{tab:seed-audit} reports every test-set concordance as stored in \texttt{results/metabric\_rfs\_benchmark.json}. Repeated splits of the same cohort are not five external replications; their spread measures sensitivity to this split protocol, not population transfer. The fold-specific percentile intervals below are the original bootstrap outputs, not confidence bounds on a five-seed average.
\begin{table}[htbp]\centering\small
\begin{tabular}{rccccc}\hline
Seed & Clinical Cox & Clinical+expression Cox & DeepSurv & CoxCNN & CoxGCN \\ \hline
0 & 0.6689 & 0.6750 & 0.6498 & 0.6632 & 0.6184 \\
1 & 0.6701 & 0.6726 & 0.6567 & 0.6632 & 0.6470 \\
2 & 0.6665 & 0.6843 & 0.6424 & 0.6574 & 0.6516 \\
3 & 0.6684 & 0.6834 & 0.6442 & 0.6549 & 0.6300 \\
4 & 0.6647 & 0.6852 & 0.6482 & 0.6604 & 0.6595 \\
\hline\end{tabular}
\caption{All five held-out seed results for the original METABRIC experiment. These are repeated partitions of one cohort.}\label{tab:seed-audit}
\end{table}
The full Cox model exceeds the clinical-only model in all five partitions, but the mean difference is small relative to the uncertainty of any individual test set. No deep architecture exceeds the full Cox mean. The cohort-level result cannot be used as a replacement for independent verification of a new signature.
\begin{table}[htbp]\centering\small
\begin{tabular}{rccc}\hline
Seed & Clinical+expression Cox & 95\% bootstrap CI & Test events \\ \hline
0 & 0.6750 & (0.6400, 0.7144) & 200 \\
1 & 0.6726 & (0.6359, 0.7081) & 200 \\
2 & 0.6843 & (0.6471, 0.7206) & 200 \\
3 & 0.6834 & (0.6484, 0.7173) & 200 \\
4 & 0.6852 & (0.6495, 0.7250) & 200 \\
\hline\end{tabular}
\caption{Individual holdout bootstrap intervals for the full Cox benchmark, retained without pooling them into a spurious external-replication interval.}\label{tab:seed-ci}
\end{table}
\subsection{External cohort accounting and comparator limits}
Table~\ref{tab:transport-audit} unpacks the five external comparisons. The frozen expression coefficients were transported without fitting to the evaluation outcomes, but within-cohort expression standardization and probe mapping remain consequential preprocessing. Point-estimate ordering is reported without a claim that overlapping intervals establish superiority. The cohorts have different event definitions, treatment settings and mapped-gene counts; pooling their concordances as if they were one common endpoint would conceal those differences.
\begin{table}[htbp]\centering\footnotesize
\begin{tabular}{lrrllp{1.25in}}\hline
Cohort & Patients & Events & Frozen panel & Comparator & Interpretation \\ \hline
GSE7390 & 196 & 89 & 0.6009 & GGI 0.5366; Veridex-76 0.5757 & RFS point-estimate beat \\
GSE25066 & 508 & 111 & 0.6400 & GGI class 0.6025 & DRFS; chemotherapy cohort \\
GSE11121 & 200 & 46 & 0.6979 & Grade 0.6294 & DMFS; node-negative \\
GSE2990 & 187 & 67 & 0.6559 & GGI continuous 0.6651 & RFS; comparator ahead \\
GSE20685 & 327 & 83 & 0.6264 & Nodal stage 0.6996 & Metastasis; comparator ahead \\
\hline\end{tabular}
\caption{Cohort-specific endpoint and comparator audit from the committed transport records. Sample numbers for the calibration audit can differ by a filtering step; the analyses are not interchangeable.}\label{tab:transport-audit}
\end{table}
\par In GSE25066 the DLDA-30 metadata direction failed an internal check against the observed outcome and was excluded; the chemotherapy-sensitivity comparator (0.6422) is effectively tied with the panel. The GSE2990 GGI score is continuous rather than a two-level published category. GSE20685 nodal stage wins by a substantial point-estimate margin, so the expression score is not a general substitute for measured clinical risk.
\subsection{Absolute-risk transfer audit}
Discrimination does not establish calibration. The 60-month METABRIC baseline survival was applied without cohort-specific refitting, and Table~\ref{tab:calibration-audit} gives the stored slope estimates. Only GSE7390 and GSE2990 have a comparable recurrence-free endpoint; the other rows are descriptive stress tests rather than absolute-risk validations. Cohort-wise standardization further limits interpretation of transported absolute risk.
\begin{table}[htbp]\centering\small
\begin{tabular}{lrlccc}\hline
Cohort & $n$ & RFS comparable? & Slope (95\% CI) & Predicted 5-y event & KM observed \\ \hline
GSE7390 & 198 & yes & 0.464 (0.077, 0.852) & 0.264 & 0.286 \\
GSE11121 & 200 & no & 0.973 (0.482, 1.463) & 0.261 & 0.146 \\
GSE2990 & 187 & yes & 1.030 (0.408, 1.652) & 0.256 & 0.248 \\
GSE20685 & 327 & no & 0.665 (0.257, 1.073) & 0.263 & 0.140 \\
GSE25066 & 507 & no & 0.687 (0.352, 1.022) & 0.264 & 0.274 \\
\hline\end{tabular}
\caption{Descriptive five-year calibration audit from \texttt{results/transport\_calibration.json}; non-RFS endpoints cannot validate RFS absolute probability.}\label{tab:calibration-audit}
\end{table}
\par GSE7390 has a slope of 0.464 (95\% CI 0.077--0.852); its overall mean predicted five-year event probability is 0.264 versus Kaplan--Meier estimate 0.286. Mean agreement conceals a shallow slope and non-monotone grouped observations. GSE2990 has slope 1.030 with a wide interval (0.408--1.652). Neither outcome justifies individual-patient risk prediction.
\subsection{Grouped calibration diagnostics for comparable endpoints}
The two recurrence-free cohorts allow descriptive quartile comparisons at 60 months. Table~\ref{tab:calibration-quartiles} retains all grouped predictions and Kaplan--Meier estimates; risk-set counts expose censoring at the target horizon. The groups are within-cohort quantiles, not common clinical cutoffs. They cannot validate the score for individual decisions or be mixed with the three differently defined endpoints.
\begin{table}[htbp]\centering\small
\begin{tabular}{llrrrr}\hline
Cohort & Quantile & $n$ & At risk at 5 y & Mean predicted & KM observed \\ \hline
GSE7390 & Q1 & 50 & 43 & 0.145 & 0.122 \\
GSE7390 & Q2 & 49 & 32 & 0.211 & 0.291 \\
GSE7390 & Q3 & 49 & 29 & 0.278 & 0.388 \\
GSE7390 & Q4 & 50 & 31 & 0.420 & 0.342 \\
GSE2990 & Q1 & 47 & 33 & 0.163 & 0.077 \\
GSE2990 & Q2 & 46 & 32 & 0.220 & 0.236 \\
GSE2990 & Q3 & 47 & 27 & 0.268 & 0.229 \\
GSE2990 & Q4 & 47 & 24 & 0.373 & 0.430 \\
\hline\end{tabular}
\caption{All recorded calibration quartiles for the two recurrence-free external cohorts. Five-year risk-set counts can be small; compare groups cautiously.}\label{tab:calibration-quartiles}
\end{table}
In GSE7390, the highest-risk quartile has lower observed recurrence at five years than the third quartile even though its predicted risk is higher, consistent with the slope warning. In GSE2990 the top quartile separates more clearly, but the middle quartiles are close. These observed curves include censoring uncertainty not displayed as confidence bands here; this descriptive table cannot substitute for prespecified calibration tests on an independently matched treatment cohort.
\subsection{Leave-one-cohort-out stability falsification}
The stability experiment trains its panel-selection rule on METABRIC and four external cohorts while withholding the fifth external cohort from gene selection. A gene must match the METABRIC coefficient direction with a bootstrap interval excluding zero in all four training cohorts. The reduced score masks unstable METABRIC coefficients; it is not re-trained on the holdout. Table~\ref{tab:stability-audit} is the complete five-fold outcome, including missing-by-design reduced scores where no gene passed the rule.
\begin{table}[htbp]\centering\small
\begin{tabular}{lrrrl}\hline
Held-out cohort & Stable genes & Full $C$ & Reduced $C$ & Selected genes \\ \hline
GSE7390 & 3 & 0.5993 & 0.5772 & CCNB1, CDC20, MELK \\
GSE11121 & 0 & 0.6963 & undefined & none \\
GSE2990 & 0 & 0.6574 & undefined & none \\
GSE20685 & 2 & 0.6264 & 0.5706 & RRM2, UBE2C \\
GSE25066 & 0 & 0.6441 & undefined & none \\
\hline\end{tabular}
\caption{Locked stability gate. A zero-gene fold has no reduced score; this is gate failure, not a zero-valued concordance.}\label{tab:stability-audit}
\end{table}
\par Three zero-gene folds make the predeclared five-fold mean delta and the planned bootstrap and random-size-matched gene-label permutation comparisons undefined. In the remaining folds, the reduced score also performs worse. Computing a mean over only those two folds would change the target of inference after seeing outcomes. The result rejects this stable-subpanel route to a biological discovery; it does not erase the separate full-panel comparator audit.
\subsection{Biology-versus-measurement boundary}
The technical audit counted annotated GPL96 and GPL570 probes for the panel genes and compared those counts with fold-stability frequency. Spearman correlations were not significant: GPL96 $\rho=-0.127$, $p=0.296$; GPL570 $\rho=-0.181$, $p=0.134$; mapped-cohort count $\rho=0.092$, $p=0.447$. Yet no gene was stable in three folds, so binary technical leave-one-gene-out AUROC is undefined. A non-significant correlation cannot rule out batch or probe bias, especially after within-cohort z-scoring has discarded raw intensity and variance information. The separately locked Hallmark coherence test likewise cannot run past its zero-gene-fold gate. Neither test licenses pathway enrichment or a causal mechanism claim.
\subsection{Fitted-model coefficient audit (bounded gene-level readout)}
The deployed risk model is the penalized Cox fit on 1,580 METABRIC patients with 17 clinical covariates and the 70-gene expression panel (87 features; held-out $n=395$, $C=0.6891$, \texttt{results/risk\_model.json}). Table~\ref{tab:coef-audit} lists the twelve largest risk-increasing and twelve largest risk-decreasing standardized coefficients. Three caveats bound every reading: the fit is penalized, so correlated features split and shrink weight; expression inputs are z-scored, so a coefficient compares one-SD expression shifts, not presence or absence of a gene; and coefficients are associational. A counterintuitive sign on any single gene is not evidence about mechanism and is not used as one anywhere in this paper.
\begin{table}[htbp]\centering\footnotesize
\begin{tabular}{lr|lr}\hline
\multicolumn{2}{c|}{Risk-increasing} & \multicolumn{2}{c}{Risk-decreasing} \\ 
Feature & $\hat\beta$ & Feature & $\hat\beta$ \\ \hline
LYMPH\_NODES\_EXAMINED\_POSITIVE & +0.2072 & HORMONE\_THERAPY=YES & -0.0659 \\
NPI & +0.1515 & BCL2L1 & -0.0681 \\
TUMOR\_SIZE & +0.1244 & ACTR3B & -0.0682 \\
UBE2C & +0.1229 & PGR & -0.0748 \\
CCNB1 & +0.1141 & MAPT & -0.0777 \\
MELK & +0.1134 & CLAUDIN\_SUBTYPE=claudin-low & -0.0803 \\
CLAUDIN\_SUBTYPE=LumB & +0.1097 & GATA3 & -0.0829 \\
FOXA1 & +0.1018 & RADIO\_THERAPY=YES & -0.0949 \\
CXXC5 & +0.0846 & ATM & -0.0966 \\
MIA & +0.0828 & ESR1 & -0.1001 \\
RB1 & +0.0823 & NDC80 & -0.1042 \\
VEGFA & +0.0789 & CHEK2 & -0.1448 \\
\hline\end{tabular}
\caption{Largest standardized coefficients of the committed risk model. Signs and magnitudes are exactly as stored; no feature was re-estimated for this table.}\label{tab:coef-audit}
\end{table}
The pattern is consistent with the rest of the paper rather than a new finding: clinical burden variables (positive lymph-node count, NPI, tumor size) dominate the risk-increasing side, joined by mitotic-panel genes (UBE2C, CCNB1, MELK), while hormone-receptor-pathway expression (ESR1, PGR, GATA3, MAPT) and recorded hormone and radiation treatment sit on the risk-decreasing side. The treatment indicators encode assignment, not benefit, and their signs reflect confounding by indication; they are reported, not interpreted as effects.
\subsection{Five-year Brier score and decision-curve audit}
Discrimination and calibration slopes do not say whether acting on the score helps. Using the frozen panel and the transported 60-month baseline (\texttt{run\_brier\_dca.py}, \texttt{results/brier\_dca.json}; design locked in the same commit as first run), Table~\ref{tab:brier-audit} reports the IPCW Brier score at 60 months against a null model that assigns every patient the cohort Kaplan--Meier five-year event probability. Inverse-probability-of-censoring weights follow the standard time-point form; patients censored before 60 months contribute through the censoring model only. The METABRIC row is apparent performance on the training cohort and is not validation.
\begin{table}[htbp]\centering\small
\begin{tabular}{llrrrr}\hline
Cohort & Role & $n$ & Brier (IPCW) & Null Brier & Skill \\ \hline
METABRIC\_apparent & apparent & 1975 & 0.1718 & 0.1910 & 0.100 \\
GSE7390 & external & 198 & 0.1988 & 0.2042 & 0.026 \\
GSE2990 & external & 187 & 0.1731 & 0.1867 & 0.072 \\
\hline\end{tabular}
\caption{IPCW Brier scores at 60 months. Skill is $1 - B/B_{\mathrm{null}}$; values near zero mean the score barely beats assigning everyone the cohort average.}\label{tab:brier-audit}
\end{table}
On the two RFS-comparable external cohorts the skill over the null is small: 0.026 in GSE7390 and 0.072 in GSE2990. The score adds a little absolute-risk information beyond the cohort average, not enough to support individual treatment decisions. Decision-curve analysis over threshold probabilities 0.05--0.60 (same IPCW weights; Table~\ref{tab:dca-audit}) shows net benefit above treat-all mainly in the 0.15--0.25 threshold band; at low thresholds the model and treat-all are nearly identical, and at 0.30 and above treat-all goes negative while the model stays positive but small.
\begin{table}[htbp]\centering\small
\begin{tabular}{lrrrr}\hline
Cohort & NB @ 0.15 & NB @ 0.20 & NB @ 0.25 & NB @ 0.30 \\ \hline
METABRIC\_apparent & 0.134 & 0.109 & 0.084 & 0.069 \\
GSE7390 & 0.172 & 0.125 & 0.059 & 0.018 \\
GSE2990 & 0.119 & 0.091 & 0.062 & 0.040 \\
\hline\end{tabular}
\caption{Decision-curve net benefit at four thresholds (treat-none is zero by construction). Descriptive audit only; thresholds for care are a clinical, not statistical, decision.}\label{tab:dca-audit}
\end{table}
These curves are descriptive: they use a transported baseline without cohort refit, assume independent censoring given time, and cannot recommend a clinical threshold. They bound, rather than promote, the score's decision value.
\subsection{Competing-risks descriptive check (cause-specific, not Fine--Gray)}
Recurrence is not the only way patients leave the risk set. Coding three states from the committed clinical records (\texttt{run\_competing\_risks.py}, \texttt{results/competing\_risks.json}; design locked in the same commit as first run): 800 recurrences, 431 deaths without recurrence, 744 censored. Death before recurrence is therefore a material competing event (about one death for every two recurrences). Cause-specific penalized Cox fits on the same clinical$+$expression features give apparent concordance 0.709 for recurrence and 0.792 for death without recurrence (in-sample, so optimistic; the validated recurrence figure remains the 0.680 cross-validated benchmark). Cause-specific models censor the other event type and answer a different question than the Fine--Gray subdistribution model, which remains not done here: no validated implementation of its time-varying censoring weights exists in this codebase, and a substituted model under the Fine--Gray name would be mislabeled. The descriptive implication stands: absolute recurrence-risk estimates that ignore the 431 competing deaths will overstate five-year recurrence probability, which bounds every calibration number in this paper.
\subsection{Gene-family-blocked cross-validation of the 70-gene panel}
Which functional families carry the panel's discrimination? The audit was preregistered before any family arm was fit; the locked design is \texttt{docs/PREREG\_FAMILY\_BLOCKED\_CV\_}\allowbreak\texttt{20260928.md}: eight a-priori families covering all 70 genes exactly once, penalized Cox (penalizer 0.05) with the four declared clinical covariates in every arm, five-fold patient cross-validation (seed 42, pooled out-of-fold concordance) plus the fixed seed-42 holdout, and 100 size-matched random gene-set nulls per family (per-draw seeds 20260928+i; \texttt{run\_family\_blocked\_cv.py}, \texttt{results/family\_blocked\_cv.json}). Baseline pooled-CV concordance 0.6550, holdout 0.6464.
\begin{table}[htbp]\centering\footnotesize
\begin{tabular}{lrccccc}\hline
Family (genes dropped) & $n$ & CV5 $C$ & $\Delta$CV5 & Holdout $C$ & $\Delta$holdout & null $\Delta$ p95 \\ \hline
basal\_myoepithelial & 10 & 0.6618 & -0.0068 & 0.6498 & -0.0034 & +0.0060 \\
cellcycle\_apoptosis\_signaling & 5 & 0.6598 & -0.0048 & 0.6458 & +0.0006 & +0.0042 \\
dna\_repair\_genome\_stability & 8 & 0.6576 & -0.0026 & 0.6507 & -0.0043 & +0.0052 \\
er\_luminal & 14 & 0.6536 & +0.0014 & 0.6426 & +0.0038 & +0.0106 \\
her2\_amplicon & 2 & 0.6552 & -0.0002 & 0.6482 & -0.0018 & +0.0034 \\
hypoxia\_adhesion\_detox\_other & 6 & 0.6542 & +0.0008 & 0.6449 & +0.0015 & +0.0060 \\
pi3k\_akt\_mtor & 4 & 0.6541 & +0.0009 & 0.6466 & -0.0002 & +0.0053 \\
proliferation\_mitotic & 21 & 0.6454 & +0.0095 & 0.6486 & -0.0022 & +0.0075 \\
\hline\end{tabular}
\caption{Leave-one-family-out concordance. $\Delta$ is baseline minus arm, so positive values mean the family helps. Under the locked rule (CV5 drop $>0.005$ AND holdout drop above the size-matched null 95th percentile) no family is load-bearing.}\label{tab:family-blocked}
\end{table}
Two honest readings follow. First, no single family is load-bearing under the locked rule: the largest cross-validation drop is proliferation\_mitotic ($+0.0095$), but its fixed-holdout drop is $-0.0022$ (percentile 50 against its size-matched null), a split verdict that echoes the earlier proliferation negative rather than contradicting it. Second, removing basal\_myoepithelial raises pooled-CV concordance to 0.6618 and is flagged dispensable-or-harmful under the same locked rule; the effect is small, the holdout does not confirm it, and it is reported as a bound on what the panel needs, not a pruning claim. The panel's discrimination is distributed across families rather than carried by any one of them.
\subsection{Patient-level endpoint harmonization before pooling}
The endpoint boundary documented above forbids pooling unlike endpoints; this audit executes the patient-level harmonization itself. It was preregistered before any harmonized endpoint was derived (\texttt{docs/PREREG\_ENDPOINT\_HARMONIZATION\_}\allowbreak\texttt{20260928.md}; raw per-patient clinical pulls and their SHA-256 manifest committed before any outcome was computed). Two strata are locked: HARM-RFS (event = any recurrence, locoregional or distant) and HARM-DMFS (event = distant metastasis or distant relapse only). A cohort enters a stratum only when its raw fields express that stratum's event; there is no imputation and no proxy mapping across endpoint types. Pooling happens only within a stratum (patient-level bootstrap stratified by cohort, 500 replicates, seed 0); the two strata are never pooled together or compared as if they were one endpoint. The model is the committed transport model - penalized Cox (penalizer 0.05) on the METABRIC 70-gene panel, fit once, with no re-tuning after any harmonized outcome was seen. The preregistration locks no performance gate: this is a methodological audit, and every number below is descriptive.
\begin{table}[htbp]\centering\footnotesize
\begin{tabular}{llrrrc}\hline
Stratum & Cohort & Patients & Events & Mapped genes & $C$ (95\% CI) \\ \hline
HARM-RFS & GSE2990 & 187 & 67 & 63 & 0.6559 (0.5778, 0.7313) \\
HARM-DMFS & GSE2990 & 179 & 40 & 63 & 0.6648 (0.5685, 0.7467) \\
HARM-RFS & GSE7390 & 198 & 91 & 63 & 0.5997 (0.5360, 0.6513) \\
HARM-DMFS & GSE7390 & 198 & 62 & 63 & 0.6352 (0.5711, 0.6954) \\
HARM-DMFS & GSE11121 & 200 & 46 & 63 & 0.6979 (0.6088, 0.7771) \\
HARM-DMFS & GSE20685 & 327 & 83 & 69 & 0.6270 (0.5638, 0.6907) \\
HARM-RFS & GSE20685 & 307 & 83 & 69 & 0.6033 (0.5397, 0.6804) \\
HARM-DMFS & GSE25066 & 508 & 111 & 63 & 0.6400 (0.5912, 0.6821) \\
\hline
HARM-RFS & \textit{pooled} (3 cohorts) & & & & 0.6165 (0.5725, 0.6577) \\
\hline
HARM-DMFS & \textit{pooled} (5 cohorts) & & & & 0.6477 (0.6191, 0.6755) \\
\hline\end{tabular}
\caption{Harmonized transport concordance by stratum. Within a stratum every cohort carries the same patient-level event definition; pooled rows are cohort-stratified bootstrap intervals, not cross-endpoint averages. GSE11121 and GSE25066 have no any-recurrence field in their committed raw pulls and enter HARM-DMFS only; GSE25066 contributes distant-relapse (DRFS) fields, which the locked definition treats as the distant-event endpoint.}\label{tab:harm-cindex}
\end{table}
Where a harmonized definition coincides with a committed transport definition the frozen pipeline reproduces the committed value exactly (GSE2990: HARM-RFS 0.6559, HARM-DMFS 0.6648), a determinism check rather than a new result. The pooled HARM-DMFS concordance 0.6477 (0.6191, 0.6755) sits above the pooled HARM-RFS concordance 0.6165 (0.5725, 0.6577), but the two strata draw on different cohort sets and their intervals overlap; the preregistration forbids reading this as an endpoint comparison. Mapped-gene counts differ by platform (63 of 70 on GPL96 cohorts, 69 of 70 on GPL570), so cross-cohort contrasts also carry a platform-coverage difference. Per-cohort intervals overlap heavily; no cohort ordering is claimed.
\begin{table}[htbp]\centering\footnotesize
\begin{tabular}{llp{2.6in}}\hline
Cohort & Stratum & Derivation from committed raw pull \\ \hline
GSE2990 & HARM-RFS & event.rfs/time.rfs verbatim (years; unit-free C-index) \\
GSE2990 & HARM-DMFS & event.dmfs/time.dmfs verbatim (years) \\
GSE7390 & HARM-RFS & e.rfs/t.rfs verbatim (GEO 'key: value' characteristics cells, key-exact, prefix stripped at parse) \\
GSE7390 & HARM-DMFS & e.dmfs/t.dmfs verbatim (same key-verbatim extraction) \\
GSE11121 & HARM-DMFS & t.dmfs/e.dmfs verbatim (months; key-verbatim extraction as GSE7390) \\
GSE11121 & HARM-RFS & EXCLUDED: no any-recurrence field in committed raw pull \\
GSE20685 & HARM-DMFS & event\_metastasis + follow\_up\_duration (years) verbatim (key-verbatim extraction) \\
GSE20685 & HARM-RFS & any recurrence = regional\_relapse OR event\_metastasis, key-verbatim (20 of 327 pulled patients carry regional\_relapse: NA, no usable value - excluded from this stratum only, leaving n=307); same follow-up field \\
GSE25066 & HARM-DMFS & drfs\_1\_event\_0\_censored + drfs\_even\_time\_years verbatim (distant relapse; key-verbatim, 198 of 508 rows recovered from drifted columns) \\
GSE25066 & HARM-RFS & EXCLUDED: no any-recurrence field in committed raw pull (neoadjuvant pCR cohort) \\
\hline\end{tabular}
\caption{Derivation audit: every stratum assignment traces to verbatim fields in the committed raw pulls. Exclusions are recorded with reasons; 20 of 327 GSE20685 patients carry \texttt{regional\_relapse:\ NA} and leave the HARM-RFS stratum only. TCGA stays out of both strata: its committed analysis is a case-control early-recurrence comparison, not a patient-level survival transport.}\label{tab:harm-deriv}
\end{table}
